% SGA 3, Expose VII_B, English reconstruction -- Loop 1.
% Scope: Section 5.1 opening, proposition, and proof reduction in full
% on authority local pages 82--83; stops immediately before Section 5.1.1.
% Combined-reader pages: 595--596. Re-edition printed pages: 584--585.
% Editorial notes in scope: 143--144.
% Authority: Polo--Gille Exp7B-13oct24.pdf,
% SHA-256 D303690E0195654BBE7C3DAA5D89EA23E24CD709A3575217234AC19AE8836694.
% Jacob Reinhold's CC BY 4.0 Markdown was consulted as comparison lineage.
% No genuine diagram occurs in this component; exact sequences are native LaTeX.

\subsection*{5.1}
\phantomsection
\label{sga3:VIIB:5.1}

设 \(k\) 为伪紧环。\footnote{%
\label{sga3:VIIB:nde:143}原文在 \SGARefLink{sga3:VIIB:5.1:proposition}{5.1} 中假设 \(k\) 为域。事实上，此假设可以换成平坦性假设；编者据此修改了 \SGARefLink{sga3:VIIB:5.1}{5.1}--\SGARefLink{sga3:VIIB:5.1.5}{5.1.5}。}
设 \(G\) 为拓扑平坦的无穷小形式 \(k\)-群，\(A\) 为其仿射代数，\(B\subset A\) 为闭子代数。置
\[
  X=\operatorname{Spf}(B),
\]
令
\[
  r:G\longrightarrow X
\]
为包含映射 \(B\subset A\) 诱导的满态射。我们将确定 \(r\) 何时使 \(X\) 成为 \(G\) 对子群 \(H\) 的右商\textup{（参见 \SGARefLink{sga3:VIIB:2.4}{2.4}）}。\footnote{%
\label{sga3:VIIB:nde:144}编者将“左”换成“右”，并修改了 \SGARefLink{sga3:VIIB:5.1:proposition}{命题~5.1} 的陈述，以便更清楚地列出等价条件 \textup{\SGARefLink{sga3:vii:paragraph:5-1:equation:i}{(i)}}、\textup{\SGARefLink{sga3:vii:paragraph:5-1:equation:ii}{(ii)}} 和结论 \(\operatorname{Spf}(B)\simeq G/H\)。}

\medskip
\noindent\textbf{命题 5.1.}
\phantomsection
\label{sga3:VIIB:5.1:proposition}
{\itshape 设 \(G\) 为拓扑平坦的无穷小形式 \(k\)-群，\(A\) 为其仿射代数，\(I_A\) 为 \(A\) 的增广理想，\(B\) 为 \(A\) 的闭子代数。置
\[
  I_B=B\cap I_A,
  \qquad
  J_B=AI_B.
\]
假设 \(A\) 以及对每个 \(n\geq0\) 的
\[
  J_B^n/J_B^{n+1}
\]
都在 \(k\) 上拓扑平坦。则下列两个断言等价：
\begin{enumerate}
\item[\textup{(i)}] 对每个 \(x\in B\)，有
\[
  \Delta_A(x)-x\widehat\otimes1
  \in
  A\widehat\otimes_kJ_B.
\]

\item[\textup{\SGAInlineRefTarget{sga3:VIIB:5.1:proposition:item:ii}(ii)}] 令
\[
  \pi:A\longrightarrow A/J_B
\]
为投影，并置
\[
  \tau_1(a)=a\widehat\otimes1.
\]
则序列
\[
  B\longrightarrow A
  \mathrel{\mathop{\rightrightarrows}^{\tau_1}_{%
    (\operatorname{id}_A\widehat\otimes\pi)\Delta_A}}
  A\widehat\otimes_k(A/J_B)
  \tag{*}\label{sga3:vii:proposition:5-1:equation:star}
\]
正合。等价地，\(B\) 是由满足
\[
  \Delta_A(x)-x\widehat\otimes1
  \in
  A\widehat\otimes_kJ_B.
\]
的所有 \(x\in A\) 组成的集合。
\end{enumerate}
此时，
\[
  H=\operatorname{Spf}(A/J_B)
\]
为 \(G\) 的形式子群，且序列
\[
  G\times H
  \mathrel{\mathop{\rightrightarrows}^{\operatorname{pr}_1}_{\lambda}}
  G
  \longrightarrow
  \operatorname{Spf}(B)
  \tag{**}\label{sga3:vii:proposition:5-1:equation:double-star}
\]
正合，其中 \(\lambda\) 为 \(G\) 的乘法在 \(G\times H\) 上的限制。因此，
\[
  \operatorname{Spf}(B)\simeq G/H.
\]}

\phantomsection
\label{sga3:VIIB:5.1:source-arrow-note}
\noindent\textit{底本说明。}\enspace
在 \textup{(*)} 中，法文 PDF 将下面的箭头标为 \((\pi\widehat\otimes\operatorname{id}_A)\Delta_A\)。紧邻此序列前后的等价条件以及所声明的右商约定都要求使用 \((\operatorname{id}_A\widehat\otimes\pi)\Delta_A\)；这里采用后者。

\medskip
\noindent\textit{证明。}
置
\[
  \overline A=A/J_B,
  \qquad
  H=\operatorname{Spf}(\overline A).
\]
则 \(H\) 为 \(G\) 的形式子簇。由于 \(J_B\subset I_A\)，增广 \(\varepsilon_A\) 诱导连续的 \(k\)-代数态射
\[
  \overline\varepsilon:\overline A\longrightarrow k.
\]

若 \textup{\SGARefLink{sga3:vii:paragraph:5-1:equation:i}{(i)}} 成立，则
\[
  \Delta_A(I_B)
  \subset
  I_B\widehat\otimes1
  +
  A\widehat\otimes J_B.
\]
因此 \(\Delta_A\) 在取商后诱导对角态射
\[
  \overline\Delta:\overline A
  \longrightarrow
  \overline A\widehat\otimes_k\overline A.
\]
映射 \(\overline\Delta\) 和 \(\overline\varepsilon\) 使 \(H\) 成为 \(G\) 的形式子幺半群。由于 \(G\) 为无穷小的，\(H\) 也是如此；于是 \SGARefLink{sga3:VIIB:2.7:prop}{命题~2.7} 表明 \(H\) 为 \(G\) 的形式子群。由 \(G/H\) 的定义\textup{（参见 \SGARefLink{sga3:VIIB:2.4}{2.4}）}，断言 \textup{\SGARefLink{sga3:vii:paragraph:5-1:equation:ii}{(ii)}} 蕴含命题的最后一个断言。

显然 \textup{\SGARefLink{sga3:vii:paragraph:5-1:equation:ii}{(ii)}} 蕴含 \textup{\SGARefLink{sga3:vii:paragraph:5-1:equation:i}{(i)}}。逆命题的证明见 \SGARefLink{sga3:VIIB:5.1.1}{5.1.1}--\SGARefLink{sga3:VIIB:5.1.5}{5.1.5}。
